% PAGES: 155-156

\noindent Symmetric matrices. Real eigenvalues of symmetric matrices. Orthogonality of the
eigenvectors of symmetric matrices.
\medskip

\noindent Symmetric positive definite matrices and symmetric positive semidefinite matrices.
\medskip

\noindent Positive definiteness criteria for symmetric matrices. Sylvester's criterion.
\medskip

\noindent Diagonal dominant symmetric positive definite matrices.
\medskip

\noindent The diagonal form of symmetric matrices.
\medskip

\section{Symmetric matrices}

\begin{definition}
A symmetric matrix is a matrix with real entries that is equal to its transpose matrix.
In other words, an $n \times n$ matrix $A$ is symmetric if and only if $A = A\tp$, i.e.,
\begin{equation}
A(j,k) \;=\; A(k,j), \quad \forall\, 1 \leq j \neq k \leq n.
\end{equation}
\end{definition}

\par\medskip\noindent\textit{Examples:}\ (i) Let $M$ be an $n \times n$ square matrix. The
matrix $A = M + M\tp$ is a symmetric matrix, since $(M\tp)\tp = M$, see (1.19), and
therefore
\[
A\tp \;=\; (M+M\tp)\tp \;=\; M\tp+(M\tp)\tp \;=\; M\tp+M \;=\; A.
\]

(ii) Let $M$ be an $m \times n$ matrix. The matrix $A = M\tp M$ is a symmetric
matrix,\footnote{Moreover, $M\tp M$ is a symmetric positive semidefinite matrix; cf.
Lemma 5.2.} since
\[
A\tp \;=\; (M\tp M)\tp \;=\; M\tp(M\tp)\tp \;=\; M\tp M \;=\; A.
\]
\hfill$\square$
\par\medskip

Recall from Theorem 4.2 that an $n \times n$ matrix with real entries has $n$
eigenvalues, counted with their multiplicities, but these eigenvalues may be complex
numbers. However, \textit{all the eigenvalues of a symmetric matrix with real entries
are real numbers}; see Theorem 5.1.

Moreover, recall from Lemma 4.2 that eigenvectors corresponding to different
eigenvalues are linearly independent. For symmetric matrices, such eigenvectors are
orthogonal; see Theorem 5.3. To state these results formally and prove them, we begin
by introducing the concepts of inner product and orthogonality.

\begin{definition}
Let $u = (u_i)_{i=1:n}$ and $v = (v_i)_{i=1:n}$ be column vectors of size $n$, with
$u_i, v_i \in \R$, for $i = 1:n$. The inner product of $u$ and $v$ is
\begin{equation}
(u,v) \;=\; v\tp u,
\end{equation}
which can be written explicitly as
\begin{equation}
(u,v) \;=\; u_1v_1 + u_2v_2 + \dots + u_nv_n \;=\; \sum_{i=1}^{n} u_iv_i.
\end{equation}
\end{definition}

Note that $v\tp u = \sum_{i=1}^{n} u_iv_i = u\tp v$, see (1.5), and therefore
\begin{equation}
(u,v) \;=\; v\tp u \;=\; u\tp v \;=\; (v,u), \quad \forall\, u,v \in \R^n.
\end{equation}

The inner product $(\cdot,\cdot)$ from (5.2) is bilinear; in particular, note that
\begin{align}
(u_1+u_2,v) &\;=\; (u_1,v)+(u_2,v), \quad \forall\, u_1,u_2,v \in \R^n;\\
(u,v_1+v_2) &\;=\; (u,v_1)+(u,v_2), \quad \forall\, u,v_1,v_2 \in \R^n;\\
(cu,v) &\;=\; (u,cv) \;=\; c(u,v), \quad \forall\, u,v \in \R^n,\ c \in \R.
\end{align}

\begin{definition}
The norm of a vector $v \in \R^n$ is defined as $\|v\| = \sqrt{(v,v)} = \sqrt{v\tp v}$.
If $v = (v_i)_{i=1:n}$, then
\begin{equation}
\|v\|^2 \;=\; (v,v) \;=\; v\tp v \;=\; \sum_{i=1}^{n} v_i^2.
\end{equation}
\end{definition}

The inner product from (5.2) and the norm from (5.8) are also called the Euclidean
inner product and the Euclidean norm of a vector, respectively.

\begin{theorem}
Any eigenvalue of a symmetric matrix is a real number.
\end{theorem}

\begin{proof}
An elegant proof of Theorem 5.1 involves an extension of the Euclidean inner product
(5.2) to complex numbers and can be found in section 10.4.
\end{proof}

\begin{lemma}
Let $A$ be an $m \times n$ matrix, and let $B$ be an $n \times m$ matrix. Let $u$ and
$v$ be column vectors of size $n$ and $m$, respectively. Then,
\begin{align}
(Au,v) &\;=\; (u,A\tp v);\\
(u,Bv) &\;=\; (B\tp u,v).
\end{align}
\end{lemma}

\begin{proof}
Recall from (1.21) that $(Bv)\tp = v\tp B\tp$ and $(A\tp v)\tp = v\tp(A\tp)\tp = v\tp A$,
since $(A\tp)\tp = A$; cf. (1.19). Then, using (5.2), we find that
\begin{align*}
(Au,v) &\;=\; v\tp Au \;=\; (v\tp A)u \;=\; (A\tp v)\tp u \;=\; (u,A\tp v);\\
(u,Bv) &\;=\; (Bv)\tp u \;=\; v\tp B\tp u \;=\; v\tp(B\tp u) \;=\; (B\tp u,v).
\end{align*}
\end{proof}
